\chapter{Conclusões e Trabalhos Futuros}
\label{cap:conclusao}

\section{Conclusões}
\label{sec:conclusoes}

Neste Trabalho de Conclusão de Curso, desenvolveu-se e validou-se um \textit{framework} preditivo computacional em linguagem Python voltado ao diagnóstico e classificação multiclasse de falhas mecânicas em rolamentos de elementos rolantes, com aplicação direcionada às caixas de mancal de rodeiros ferroviários. A fundamentação do estudo ancorou-se no banco de dados de referência internacional da \textit{Case Western Reserve University} (CWRU), contemplando séries temporais acelerométricas adquiridas a \SI{12}{\kilo\hertz} em quatro condições operacionais (saudável, falha na pista interna, falha na pista externa e falha em elementos rolantes) sob diferentes severidades de defeito e patamares de carregamento.

O cumprimento sistemático dos objetivos específicos delineados permitiu estabelecer as seguintes conclusões científicas e de engenharia:
\begin{enumerate}
    \item \textbf{Eficácia da Engenharia de Atributos Híbrida:} A combinação sinérgica entre 8 descritores estocásticos no domínio do tempo (com destaque para o Desvio Padrão, amplitude Pico a Pico e valor eficaz RMS) e 8 características no domínio da frequência (notadamente o Centroide Espectral e as energias nas frequências características de defeito BPFO, BPFI e BSF extraídas do espectro de envelope via Transformada de Hilbert com remoção de offset DC) mostrou-se decisiva. Conforme demonstrado na validação cruzada, modelos treinados exclusivamente com estatísticas temporais apresentaram desempenhos inferiores e forte instabilidade entre dobras; a introdução dos parâmetros espectrais de envelope convergiu as acurácias médias para patamares superiores a \SI{99,4}{\percent}, com desvios padrão inferiores a $\pm \SI{0,50}{\percent}$;
    \item \textbf{Prevenção Rigorosa de Vazamento de Dados (Data Leakage):} A metodologia de particionamento e validação cruzada agrupada por arquivo físico de ensaio (\textit{Stratified Group K-Fold}) garantiu que nenhuma janela temporal de um ensaio de teste compartilhasse histórico ou parâmetros com a fase de treino. Essa blindagem metodológica assegurou que as métricas reportadas reflitam a real capacidade de generalização dos algoritmos frente a mancais nunca antes processados, superando limitações de estudos que utilizam divisões aleatórias ingênuas de janelas;
    \item \textbf{Desempenho dos Classificadores Supervisionados:} No conjunto de teste independente composto por 12 ensaios físicos (\num{886} janelas discretas de \SI{0,17}{\second}), o classificador \textit{Random Forest} atingiu acurácia global de \SI{99,32}{\percent} e \textit{F1-Score} macro de \num{0,9926}, enquanto o modelo \textit{Support Vector Machines} (SVM com \textit{kernel} RBF e padronização \textit{StandardScaler}) alcançou acurácia global de \SI{99,10}{\percent} e \textit{F1-Score} macro de \num{0,9900}. Ambas as arquiteturas apresentaram sensibilidade de \SI{100,00}{\percent} para a detecção de mancais saudáveis e falhas em pista externa;
    \item \textbf{Compreensão Mecatrônica dos Desvios na Classe B:} O recall ligeiramente inferior observado na detecção de defeitos em elementos rolantes (esferas) foi fisicamente interpretado à luz da cinemática planetária dos corpos rolantes, cujo movimento tridimensional não confinado pode amortecer os pulsos de impacto quando a cavidade transita fora do plano de carregamento principal em condições de baixa carga estática;
    \item \textbf{Confiabilidade Plena via Votação Majoritária por Ensaio:} A implementação da estratégia de votação por maioria temporal em nível de ensaio (\textit{file-level voting}) atingiu \textbf{\SI{100,00}{\percent}} de acerto em todos os 12 ensaios do conjunto de teste independente. Essa convergência confirma que, para fins práticos de engenharia mecatrônica ferroviária, o monitoramento contínuo em pequenos blocos de tempo elimina integralmente os erros espúrios de janelas isoladas, fornecendo um diagnóstico determinístico e isento de falsos alarmes;
    \item \textbf{Qualidade de Software e Arquitetura de Dados:} O pipeline computacional foi estruturado em módulos independentes com tipagem estrita (\textit{type hints}), contratos automatizados de validação de dados e persistência em formato colunar Apache Parquet, garantindo plena reprodutibilidade e conformidade com as melhores práticas de Engenharia de Dados.
\end{enumerate}

\section{Trabalhos Futuros}
\label{sec:trabalhos_futuros}

Como desdobramentos desta pesquisa, propõem-se as seguintes frentes de continuidade:
\begin{enumerate}
    \item \textbf{Implementação em Sistemas Embarcados de Borda (Edge AI / TinyML):} Embarcar o pipeline de extração de atributos e o modelo leve otimizado (como árvores de decisão quantizadas) em nós sensores inteligentes equipados com microcontroladores de baixo consumo de energia (ex.: arquiteturas ARM Cortex-M4/M7), permitindo o processamento local nas caixas de mancal dos vagões sem dependência contínua de telemetria em tempo real;
    \item \textbf{Expansão para Classificação Hierárquica de 10 Classes:} Estender o classificador para prever não apenas o tipo de falha mecânica, mas também quantificar diretamente o nível de severidade do defeito (\SI{0.007}{in}, \SI{0.014}{in} e \SI{0.021}{in}), viabilizando a estimativa da Vida Útil Remanescente (\textit{Remaining Useful Life} -- RUL) do componente;
    \item \textbf{Fusão Sensorial Heterogênea:} Integrar ao vetor de características atributos derivados de sensores térmicos pontuais de alta frequência, acelerômetros triaxiais instalados no corpo da caixa de mancal e transdutores de Emissão Acústica (EA), investigando técnicas de fusão de dados em nível de atributos ou em nível de decisão;
    \item \textbf{Validação em Dados Reais de Malha Ferroviária de Carga:} Coletar e aplicar a metodologia a sinais acelerométricos registrados em condições operacionais de campo de ferrovias pesadas (\textit{heavy haul}), submetendo os modelos aos efeitos dinâmicos de velocidade angular variável, acelerações de tração/frenagem e irregularidades geométricas da via permanente;
    \item \textbf{Investigação de Redes Neurais Profundas (Deep Learning):} Explorar abordagens de aprendizado profundo de ponta a ponta (\textit{end-to-end}), como Redes Neurais Convolucionais 1D (CNN-1D) e redes recorrentes com mecanismos de atenção aplicadas sobre espectrogramas de tempo-frequência (Transformada de Wavelet Contínua -- CWT), confrontando seu custo computacional com o dos modelos clássicos avaliados neste estudo.
\end{enumerate}
